\section{Related Work}
\label{sec:related}
This section relates AgentWorkShop to plant-facing industrial agents and to
production and automation abstractions for context, tool access, and
operational constraints. Table~\ref{tab:nearest} compares the integration
features explicitly described in selected architecture-level studies closest
to the evaluated task-to-device path. The comparison is focused rather than
exhaustive; a feature marked as not specified was not described in the cited
study and is not inferred to be absent.

\subsection{Industrial Agents and Plant Integration}
Industrial-agent systems, including large language model (LLM)-based systems,
span end-to-end automation control, supervisory
services, programmable logic controller (PLC) interfaces, manufacturing
coordination, and operator-facing machine access. Xia et al. present an
LLM agent framework for end-to-end industrial automation control, using
structured prompts and event-driven context for production planning and
operation \cite{xia2025control}. Huang et al. present a multimodal
architecture whose stated focus is automated control generation for
industrial cyber-physical systems \cite{huang2026made}. Biskupovic et al.
study agent-based supervision for service-oriented industrial systems
\cite{biskupovic2025supervision};
Agents4PLC and InstructMPC address PLC code generation and human--LLM
context-aware control, respectively
\cite{liu2026agents4plc,wu2025instructmpc}. Manufacturing studies consider
dynamic resource-capability exploration \cite{lim2025llmenabled}, multi-agent
shop-floor coordination \cite{zhao2026shopfloor}, and digital-twin-supported
fault handling \cite{gill2025llmdt}. Natural-language machine access
\cite{hofmann2025nlc} and real-time LLM control of a robot-operated serial
line \cite{waseem2026pretrained} extend this work to operator interfaces and
robotic-line control. Together, these studies address different operational
contexts, interfaces, and allocations of responsibility between agents,
operators, and controllers.

The closest architecture-level comparison is Xia et al., who connect cognitive,
digital-twin, and physical-automation layers and organize twin information as
state snapshots, event messages, and plan sequences; their case studies
address robotic automation and process simulation
\cite{xia2025architecture}. Other nearby designs expose factory-module
operations through MCP with state-based tool filtering, plan constrained
manufacturing steps, or compose semantic recipes from OPC UA skills
\cite{koehle2026llmdriven,silva2025mcpmanufacturing,koehle2025semanticrecipes}.
These studies describe important parts of plant-facing integration. The
specific design focus here is a joinable execution record across a registered
task, its per-execution production context, approval, verified route receipt,
and a later observation selected by active run. R20 checks this record path
across two signal/protocol mappings; it does not measure the path's effect on
process quality or its generality beyond the tested configuration.
\begin{table*}[!t]
\centering
\caption{Feature-level comparison of selected plant-facing integration
studies. ``Not specified'' means that the cited study does not describe the
feature; it is not evidence of absence. The table compares stated designs and
evaluation evidence, not performance.}
\label{tab:nearest}
\small
\setlength{\tabcolsep}{3pt}
\begin{tabularx}{\textwidth}{@{}p{1.8cm}p{3.0cm}p{2.4cm}p{2.8cm}X@{}}
\toprule
System & Task and production context & Per-call operation checks & Device-facing route & Reported evaluation and record link \\
\midrule
Xia et al. \cite{xia2025architecture} & Cognitive, twin, and physical layers; snapshots, events, and plans & Not specified & Digital-twin/physical-automation integration & Robotic automation and process-simulation case studies \\
K{\"o}hle et al. \cite{koehle2026llmdriven} & Factory-module agents with injected live state & State-based tool filtering & Per-module MCP exposes OPC UA methods; MQTT coordination & Nine simulated production challenges \\
Vieira da Silva et al. \cite{silva2025mcpmanufacturing} & LLM plans manufacturing steps under constraints & Not specified & MCP-exposed manufacturing resource functions & Laboratory multi-step process with constraint handling \\
K{\"o}hle et al. \cite{koehle2025semanticrecipes} & AI-assisted recipes composed from OPC UA skills & Not specified & OPC UA skills and virtual commissioning & Unity assembly-workstation proof of concept \\
\midrule
This work & Collaborative task with per-turn line--product--recipe--run context & Scope, policy, and applicable grant/binding checked before handler selection & Semantic node resolves to configured DAQ/DCW driver & Six R20 executions pass task, approval, verified-receipt, and run-tagged-sample linkage checks across two mappings; all samples remain outside target \\
\bottomrule
\end{tabularx}
\end{table*}

The differentiating object is the record chain, not an individual protocol,
agent-reasoning method, or process-control algorithm: the implementation keeps
prompt grounding, per-call authority, device readback, and run-filtered process
observation as separate states, then links the records through available
identifiers. The cited studies cover related architectural elements, but the
selected comparison does not report this same task-to-receipt-to-observation
trace as an evaluation unit. R20 demonstrates record closure for six
executions; it does not establish a performance advantage over those systems.
Authority and command confirmation are operational services consistent with
safety-aware decision support \cite{jiang2026agentic}; they do not replace
process protection or establish operator effectiveness.

\subsection{Production Context, Agent Interfaces, and Automation}
Industrial integration also depends on abstractions that are not specific to
agent architectures. MCP standardizes tool discovery and calls with
request-scoped authorization \cite{mcp2024}, whereas
ISA-88 and ISA-95 describe recipes and production operations
\cite{isa88_2010,isa95_2025}. Run-to-run methods connect process observations
to successive production runs \cite{moyne2000r2r}. At the agent level,
AgentSpec and CaMeL address tool behavior and capability-based information
flow \cite{wang2025agentspec,debenedetti2025camel}. Together, these works
provide complementary abstractions for production context, tool interaction,
and information flow. AgentWorkShop composes such abstractions with its task
runtime, operation checks, semantic I/O, and route-specific records; this
paper evaluates selected paths through that composition rather than proposing
a replacement for the underlying standards.

Industrial security and record requirements further constrain that
architecture: FDA guidance discusses trustworthy records and audit trails
\cite{part11}, IEC~62443-3-3 specifies system-security requirements
\cite{iec62443-3-3}, and DNP3 includes select-before-operate
\cite{ieee1815dnp3}. Control-assurance methods address state evolution through
shielding, control-barrier functions, or Simplex architectures
\cite{alshiekh2018shielding,ames2017cbf,sha2001simplex}. These standards and
methods inform the boundaries within which an industrial-agent framework must
operate; they are not contributions claimed by AgentWorkShop.

The evaluation tests selected software paths and simulator scenes, not
physical-plant performance.
